\section{Introduction}

Much of active equity capital is in concentrated portfolios: technology and thematic funds, sector funds, best-ideas books and the systematic baskets on narrow universes used in structured products. These managers are measured on drawdowns. Allocators gate capital on peak-to-trough losses, performance fees are paid above a high-water mark, and structured overlays are priced from realised volatility and drawdown. The objective in this paper is therefore the Calmar ratio (compounded return over maximum drawdown), with the Sharpe ratio reported next to it.

Concentration creates three problems which the literature on diversified factor portfolios mostly avoids. The first is low breadth. By the fundamental law of active management \citep{grinold1989,clarke2002} the information ratio scales with the square root of the number of independent bets, and ten stocks give very few. The second is that one factor dominates. A market or sector factor explains most of the variance of a concentrated universe (average pairwise correlation 0.40 in the development universe below), and nearly every signal is correlated with beta or momentum. The third is camouflage. Scaling net exposure with a regime model or a trend filter can raise the Calmar ratio of a weak selection book through market beta, above all in a rising market, and the statistics then look like skill. With low breadth the selection component is also measured imprecisely, so the overlay can dominate reported performance.

There are two questions here. The practical one is whether a concentrated manager can obtain a stock-selection alpha independent of market exposure, and improve the Calmar ratio with that alpha rather than with beta. The academic one is how to validate such an alpha when the sample is too small for conventional tests to have power and many specifications have been tried.

I use sell-side target prices for both. They are available point in time for every listed large cap and their information content is well documented \citep{brav2003,asquith2005}. The research had four stages, all reported.

\para{Stage 1) Engineering solved, alpha missing:} The first-generation strategy, Regime-Timed Breadth (RTB), combines a revision-breadth composite with GJR-GARCH risk weights and a hidden-Markov regime tilt. It is built with strict point-in-time rules and reaches a development Calmar ratio of 1.77. In the decomposition the timing leg earns 11.0\% a year and the selection leg has a Calmar ratio of only 0.36. Timing regressions of the \citet{treynor1966} and \citet{henriksson1981} type find no timing ability, and a static exposure of the same average size beats the regime tilt, so most of the 1.77 is beta (Figure~\ref{fig:camouflage}).

\para{Stage 2) A new signal and a horizon-matched construction:} I classify each revision by its arrival time and call the resulting signal leader minus herd (LMH). Leaders publish on an information-event day or after at least ten sessions with no other revision on the stock, as in \citet{cooper2001}. Herders publish 3--15 sessions after an earnings cluster and in its direction, in line with career-concern herding \citep{welch2000,hong2000,clement2005}. Leader breadth has a rank IC of 0.066, herder breadth $-0.084$. LMH is the only one of 18 candidate signals with incremental power over four price factors (Fama--MacBeth $t = 2.46$). With three breadth signals it enters an equal-weight rank composite, traded in 42 daily overlapping tranches. This zero-beta book, First-Mover Flow, market-neutral (FMF-N), has a net Sharpe ratio of 0.98 in development, 1.33 in the 2020--21 holdout and 1.13 overall, and about twice the selection Calmar ratio of RTB at zero beta.

\para{Stage 3) Adversarial validation:}
\begin{itemize}
  \item \emph{Institutional comparables:} FMF-N is compared with 13 signals and 3 allocations used by institutional desks. It has the highest Sharpe and Calmar ratios of the cross-sectional books built identically, returns $+12.0\%$ in the COVID crash against $-35.5\%$ for the market, keeps an alpha of 6.1\% ($t = 2.51$) against the comparables not built from its components, and improves 16 of 18 books in a blend
  \item \emph{Concentration:} In a full leave-one-stock-out rebuild no single name carries the result
  \item \emph{Never-seen period:} The main test. Under a hashed pre-registration, the frozen specification is applied to 30 stocks over 2023-2026, a period not used for any choice (twenty new names, ten development names). Sharpe ratio 1.05, Calmar ratio 1.32, alpha $t$-statistic 2.04, beta zero 2nd of 14 cross-sectional books on both measures. LMH replicates with essentially the same information coefficient
\end{itemize}

\para{Stage 4) Attacking the result:} Section~\ref{sec:objections} takes nine objections raised during the analysis and answers each with a test chosen to break the result.
\begin{itemize}
  \item \emph{Hindsight:} A causal split into a static tilt and a timing component puts the edge in timing (Sharpe 1.04 in development, 1.33 out of sample). The static tilt towards later winners is insignificant and loses money out of sample
  \item \emph{Implementation:} One-session delay, 100\,bp borrow fee, 10\,bp costs: Sharpe ratio 1.06
  \item \emph{Research degrees of freedom:} In a space of about 300 specifications, including the quarterly target expiry an investment committee would ask about, the sign is positive in 99\% of cases the selection-free ensemble earns 0.92
  \item \emph{The new signal:} Adding leader minus herd to classic breadth raises the Sharpe ratio by 0.28 in sample ($p = 0.028$) and by a similar amount out of sample
  \item \emph{Small samples:} Bayesian posteriors under sceptical priors, a random-effects meta-analysis over three independent samples and empirical-Bayes shrinkage of the signal search give the probabilities an allocator needs in place of single $p$-values
\end{itemize}
Section~\ref{sec:generative} turns the Calmar ratio from one draw into a distribution, using among other models a hybrid of GJR-GARCH dynamics and a conditional diffusion model for the joint shocks.

On the economic side, the paper shows that the timing of analyst revisions is informative in addition to their direction, and that in a concentrated universe it gives a zero-beta alpha which blends well with the factor strategies institutions already hold. On the method side, it gives a validation protocol for small samples that have been searched many times: calibrated non-overlapping tests, placebos that keep signal persistence, deflation over the full research log \citep{bailey2014,bailey2017}, a multiverse with family-wise tests \citep{hansen2005,romano2005}, Bayesian meta-analysis across samples, pre-registration and replication. The Python package (open source with tests) reproduces every number and runs the protocol, a production book, a QIS index and generative stress tests on other universes. The protocol limits the claims too. With ten stocks the detectable information coefficient is about 0.138 or more, the development search alone is only suggestive, and momentum crashes are a real risk for the book. I report these limits with the results, and list more than fifty signals and hypotheses that failed.

Literature is in Section~\ref{sec:lit}, Section~\ref{sec:framework} formalises the three problems and the protocol data are in Section~\ref{sec:data}, RTB is analysed in Section~\ref{sec:rtb} and First-Mover Flow defined in Section~\ref{sec:fmf}. Development and holdout results are in Section~\ref{sec:results}, institutional comparables in Section~\ref{sec:toolkit}, the out-of-sample replication in Section~\ref{sec:oos}, the nine objections in Section~\ref{sec:objections} and generative path risk in Section~\ref{sec:generative}. Section~\ref{sec:discussion} covers practical use and Section~\ref{sec:conclusion} concludes.
